%============================================================
% CH2 · PATTERN 2 — SUBTREE AGGREGATION AND PUSH-DOWN: THEORY
%============================================================
\subsection{Pattern Theory}
\label{ssec:ch2-p2-theory}

% the tree T of the toolbox lesson as a pic with per-vertex styles (prefix g2p2-)
\tikzset{
    g2p2-v1/.style=gv, g2p2-v2/.style=gv, g2p2-v3/.style=gv, g2p2-v4/.style=gv,
    g2p2-v5/.style=gv, g2p2-v6/.style=gv, g2p2-v7/.style=gv, g2p2-v8/.style=gv,
    g2p2-e/.style=gtree,
    g2p2-T/.pic={
        \node[g2p2-v1] (t1) at (0,0) {1};
        \node[g2p2-v2] (t2) at (-1.8,-1) {2}; \node[g2p2-v3] (t3) at (0,-1) {3}; \node[g2p2-v4] (t4) at (1.8,-1) {4};
        \node[g2p2-v5] (t5) at (-2.5,-2) {5}; \node[g2p2-v6] (t6) at (-1.1,-2) {6}; \node[g2p2-v7] (t7) at (1.8,-2) {7};
        \node[g2p2-v8] (t8) at (1.8,-3) {8};
        \foreach \a/\b in {1/2, 1/3, 1/4, 2/5, 2/6, 4/7, 7/8} \draw[g2p2-e] (t\a) -- (t\b);
    }
}

\begin{patternbox}{Pattern 2 at a Glance}
Root the tree. Information flows in two directions along the DFS: \textbf{up}, at the moment a vertex is \emph{left} (post-order) --- every child is finished, so the vertex combines its children into a value for its whole subtree; and \textbf{down}, at the moment a vertex is \emph{entered} (pre-order) --- its parent is finished, so the vertex extends the parent's value by one step. One pass, \VE{} $= \bigO{n}$.

\begin{center}
\begin{tikzpicture}[x=1.05cm, y=0.95cm]
    \tikzset{g2p2-v1/.style=gvsrc, g2p2-v4/.style=gvcur, g2p2-v7/.style=gvdone, g2p2-v8/.style=gvdone}
    \pic {g2p2-T};
    \draw[->, thick, accentorange] (0.45,0.25) to[bend left=30] node[pos=0.5, above right, font=\scriptsize, text=accentorange] {down} (2.15,-0.6);
    \draw[->, thick, primaryblue] (2.45,-3.0) -- node[pos=0.5, right, font=\scriptsize, text=primaryblue] {up} (2.45,-1.3);
    \node[font=\scriptsize, anchor=west, align=left] at (4.0,-0.5) {\textcolor{accentorange}{\textbf{enter 4:}} $\mathrm{dep}[4] = \mathrm{dep}[1] + 1 = 1$\\[1pt]\quad the parent 1 (green) is already final};
    \node[font=\scriptsize, anchor=west, align=left] at (4.0,-2.1) {\textcolor{primaryblue}{\textbf{leave 4:}} $\mathrm{sz}[4] = 1 + \mathrm{sz}[7] = 3$\\[1pt]\quad the subtree $\{7, 8\}$ (blue) is already final};
\end{tikzpicture}
\end{center}
\begin{recurrencebox}
\centering\small
\textbf{up} (leave $v$): $\ \mathrm{agg}[v] = \mathrm{base}(v) \oplus \textstyle\bigoplus_{c} \mathrm{agg}[c]$ over the children $c$
\qquad e.g.\ $\mathrm{sz}[v] = 1 + \sum_c \mathrm{sz}[c]$\\[3pt]
the part above $v$ has $n - \mathrm{sz}[v]$ vertices\\[3pt]
\textbf{down} (enter $v$): $\ \mathrm{acc}[v] = \mathrm{acc}[\mathrm{par}(v)] \otimes \mathrm{step}(v)$
\qquad e.g.\ $\mathrm{dep}[v] = \mathrm{dep}[\mathrm{par}(v)] + 1$\\[3pt]
no recursion: one BFS gives \texttt{order}; forward over \texttt{order} $=$ down, backward $=$ up
\end{recurrencebox}
\end{patternbox}

%------------------------------------------------------------
\subsubsection*{From the previous pattern to this one}

Pattern~1 looked at the edges DFS meets that are \emph{not} tree edges. A tree has none: every edge is a tree edge, the only visited neighbour of a vertex is its parent, and the DFS from the root \emph{is} the rooted tree. So this pattern ignores edge types and uses the other thing Lesson~\ref{ssec:ch2-call-stack} showed: the two hooks. Code in the post-order hook runs when a vertex's whole subtree is done; code in the pre-order hook runs when the whole path from the root is done. The toolbox lesson (\S\ref{ssec:ch2-toolbox}) computed depth and size this way on the tree $T$; here we make that program into a pattern and remove its recursion.

Three things change compared with a DFS on a general graph:
\begin{itemize}[leftmargin=*]
    \item \textbf{No \texttt{visited} array.} Pass the parent instead: in a tree, ``$v \ne \mathrm{par}(u)$'' is the same as ``$v$ not visited''. This is the only way the input's $n - 1$ undirected edges become a rooted tree.
    \item \textbf{Values, not marks.} Each vertex stores a number (or a small record) computed from its children or from its parent.
    \item \textbf{The search order is irrelevant.} Only the tree matters, so any order in which parents come before children works --- including BFS order, or simply the labels when the input guarantees $p_i < i$.
\end{itemize}

%------------------------------------------------------------
\subsubsection*{When to recognise it (and when not)}

\begin{itemize}[leftmargin=*]
    \item ``\textbf{for every vertex}, how many / what sum / which maximum \textbf{in its subtree}'' --- post-order aggregation (Problem~\ref{ssec:ch2-white-black});
    \item ``remove a vertex / cut an edge: how big are the pieces?'' --- the pieces are the child subtrees and the part above, of size $n - \mathrm{sz}[v]$ (Problem~\ref{ssec:ch2-highest-score});
    \item ``\textbf{add to the whole subtree} of $p$'', ``a value accumulated along the path from the root'', depths, ancestors' properties --- pre-order push-down (Problem~\ref{ssec:ch2-ki});
    \item ``everything that passes through the edge above $v$'' --- traffic, cars, fuel, people: it is a function of the subtree of $v$ (E2.3);
    \item a greedy that cuts a subtree as soon as it is ``complete'' (even size, sum divisible by $k$): decide at the moment $v$ is left (E2.2, E2.5).
\end{itemize}
\textbf{Not this pattern:}
\begin{itemize}[leftmargin=*]
    \item the answer is wanted \textbf{for every possible root} (``sum of distances from every vertex'') --- that is \emph{rerooting}, a second pass that moves the root across each edge: \dpbook{19}{Rerooting};
    \item the value of a vertex is a \textbf{choice} that depends on choices in the children (take / skip a vertex, colourings, matchings, knapsack over subtrees) --- general tree DP: \dpbook{11}{DP on Trees};
    \item many \textbf{queries} or \textbf{updates on subtrees mixed in time} (``add to the subtree of $p$, then ask the sum of the subtree of $q$, then \dots'') --- flatten the tree with entry times (Pattern~3) and use a Fenwick or segment tree (Chapter~15);
    \item questions about \textbf{two vertices} (distance, lowest common ancestor) --- Chapter~14.
\end{itemize}

%------------------------------------------------------------
\subsubsection*{The algorithm and its invariant}

\begin{center}
\small
\renewcommand{\arraystretch}{1.3}
\begin{tabularx}{\linewidth}{@{}>{\bfseries}l X X@{}}
\toprule
 & \textbf{up (aggregation)} & \textbf{down (push-down)} \\
\midrule
hook & post-order: after the loop over the children & pre-order: before the loop \\
needs & every child $c$ is finished & the parent is finished \\
typical values & size, number of white vertices, sum, max, product of child sizes, a top-3 list & depth, sum along the root path, ``best path ending here'', total added by ancestors \\
iterative form & loop over \texttt{order} \textbf{backwards} & loop over \texttt{order} \textbf{forwards} \\
\bottomrule
\end{tabularx}
\end{center}

\begin{invariantbox}[When $v$ is left, every value of its subtree is final]
\textbf{Invariant (up).} At the moment the DFS leaves $v$, it has entered and left every descendant of $v$ and no other vertex since entering $v$; hence $\mathrm{agg}[c]$ is final for every child $c$, and $\mathrm{agg}[v]$ computed now is final too.

\emph{Proof by induction on the height of $v$.} A leaf has no children: $\mathrm{agg}[v] = \mathrm{base}(v)$. Otherwise the loop of $v$ calls the DFS on each child $c$ and returns only after $c$ is left; by induction $\mathrm{agg}[c]$ is final then, and the subtrees of different children are disjoint (a tree has no cross edges), so nothing is counted twice. The subtree of $v$ is $\{v\} \cup \bigcup_c \mathrm{subtree}(c)$, which is exactly what the formula combines.

\textbf{Invariant (down).} At the moment the DFS enters $v$, the value of every ancestor of $v$ is final --- they are exactly the vertices on the stack (Lesson~\ref{ssec:ch2-call-stack}), each entered before $v$.

\emph{Iterative form.} The same two statements hold for any \texttt{order} in which $\mathrm{par}(v)$ precedes $v$: going backwards, all children of $v$ appear after $v$ in \texttt{order}, so they are processed before $v$; going forwards, the parent is processed before $v$. \hfill$\blacksquare$
\end{invariantbox}

\begin{diagrambox}
\begin{tikzpicture}[x=0.9cm, y=0.85cm]
    \begin{scope}
        \pic {g2p2-T};
        \foreach \v/\s/\b/\w in {1/8/36/right, 2/3/13/left, 3/1/3/below, 4/3/19/right, 5/1/5/below, 6/1/6/below, 7/2/15/right, 8/1/8/right}
            \node[gdist, text=primaryblue, fill=white, \w=1pt of t\v] {\s\,|\,\b};
        \node[font=\small\bfseries] at (0,1.0) {up: $\mathrm{sz}\,|\,\mathrm{sub}$};
        \node[note, text width=5.2cm] at (0,-4.1) {values $\mathrm{val}(v) = v$; $\mathrm{sub}[4] = 4 + \mathrm{sub}[7] = 4 + 15 = 19$};
    \end{scope}
    \begin{scope}[xshift=7.6cm]
        \pic {g2p2-T};
        \foreach \v/\d/\p/\w in {1/0/1/right, 2/1/3/left, 3/1/4/below, 4/1/5/right, 5/2/8/below, 6/2/9/below, 7/2/12/right, 8/3/20/right}
            \node[gdist, text=accentorange, fill=white, \w=1pt of t\v] {\d\,|\,\p};
        \node[font=\small\bfseries] at (0,1.0) {down: $\mathrm{dep}\,|\,\mathrm{path}$};
        \node[note, text width=5.2cm] at (0,-4.1) {$\mathrm{path}[8] = \mathrm{path}[7] + 8 = 12 + 8 = 20 = 1 + 4 + 7 + 8$};
    \end{scope}
\end{tikzpicture}
\figcap{Both directions on $T$ with $\mathrm{val}(v) = v$ (numbers printed by the templates below). Left: subtree size and subtree sum, computed bottom-up. Right: depth and the sum along the root path, computed top-down.}
\end{diagrambox}

%------------------------------------------------------------
\subsubsection*{The part above $v$: $n - \mathrm{sz}[v]$}

Deleting a vertex $v$ splits the tree into the subtrees of its children and \emph{one} more piece: everything that is not in the subtree of $v$. That piece has $n - \mathrm{sz}[v]$ vertices (zero for the root). So after one post-order pass we know the size of every piece of every deletion, although the DFS never ``looked upwards''. Products of piece sizes (Problem~\ref{ssec:ch2-highest-score}), the largest piece (centroids, Chapter~14), and the number of paths through $v$ are all built from these numbers.

\begin{diagrambox}
\begin{tikzpicture}[x=0.9cm, y=0.85cm]
    \begin{scope}
        \tikzset{g2p2-v2/.style=gvdead, g2p2-v5/.style=gvA, g2p2-v6/.style=gvB,
                 g2p2-v1/.style=gvsel, g2p2-v3/.style=gvsel, g2p2-v4/.style=gvsel, g2p2-v7/.style=gvsel, g2p2-v8/.style=gvsel}
        \pic {g2p2-T};
        \node[font=\small\bfseries] at (0,1.0) {delete 2};
        \node[note, text width=5.6cm] at (0,-4.5) {pieces $\{5\}$, $\{6\}$ (children, size 1 each) and the purple rest of size $n - \mathrm{sz}[2] = 8 - 3 = 5$; largest piece 5};
    \end{scope}
    \begin{scope}[xshift=7.6cm]
        \tikzset{g2p2-v1/.style=gvdead, g2p2-v2/.style=gvA, g2p2-v5/.style=gvA, g2p2-v6/.style=gvA,
                 g2p2-v3/.style=gvB, g2p2-v4/.style=gvsel, g2p2-v7/.style=gvsel, g2p2-v8/.style=gvsel}
        \pic {g2p2-T};
        \node[font=\small\bfseries] at (0,1.0) {delete the root 1};
        \node[note, text width=5.6cm] at (0,-4.5) {only child pieces: $3, 1, 3$; the part above is empty ($8 - \mathrm{sz}[1] = 0$); largest piece 3 --- the smallest ``worst'' of all vertices};
    \end{scope}
\end{tikzpicture}
\figcap{Deleting a vertex: the child subtrees plus the part above. Template (4) below prints the largest piece for every $v$: $3, 5, 7, 5, 7, 7, 6, 7$.}
\end{diagrambox}

%------------------------------------------------------------
\subsubsection*{No recursion: the order array}

With $n = 2\cdot10^5$ a path-shaped tree makes a recursive DFS go $2\cdot10^5$ levels deep, and on a 2\,MB stack that crashes near $65\,000$ (Lesson~\ref{ssec:ch2-call-stack}). Neither direction needs the DFS rhythm, only ``parents before children'', so:
\begin{enumerate}[leftmargin=*]
    \item run a BFS from the root, recording \texttt{order} and \texttt{parent} (skip the parent instead of checking \texttt{visited});
    \item \textbf{forward} over \texttt{order}: push-down values;
    \item \textbf{backward} over \texttt{order}: aggregate into the parent.
\end{enumerate}
If the input is a parent array with $p_i < i$ (Codeforces loves this format), even the BFS is unnecessary: the labels $1, 2, \dots, n$ already are such an order.

\begin{diagrambox}
\begin{tikzpicture}[x=1.15cm]
    \node[font=\scriptsize\bfseries, anchor=east] at (-0.6,0) {order};
    \foreach \v [count=\i from 0] in {1,2,3,4,5,6,7,8}
        \node[cell, minimum width=9mm] (o\i) at (\i,0) {\v};
    \node[font=\scriptsize\bfseries, anchor=east] at (-0.6,-0.75) {parent};
    \foreach \p [count=\i from 0] in {--,1,1,1,2,2,4,7}
        \node[idx] at (\i,-0.75) {\p};
    \draw[->, thick, accentorange] (-0.3,1.25) -- node[above, font=\scriptsize]{forward: $\mathrm{dep}[v] = \mathrm{dep}[\mathrm{par}]+1$ (parent already done)} (7.3,1.25);
    \draw[->, thick, primaryblue] (7.3,-1.35) -- node[below, font=\scriptsize]{backward: $\mathrm{sz}[\mathrm{par}] \mathrel{+}= \mathrm{sz}[v]$ (children already done)} (-0.3,-1.35);
    \draw[dep] (o7.north) to[bend right=40] (o6.north);
    \draw[dep] (o6.north) to[bend right=40] (o3.north);
    \node[note, text width=3.2cm, anchor=west] at (8.0,-0.4) {BFS order of $T$ from 1; orange arcs: 8 hands its size to 7, then 7 to 4};
\end{tikzpicture}
\figcap{The order-array trick. BFS gives every parent before its children, so one forward and one backward loop replace both hooks of the recursive DFS.}
\end{diagrambox}

%------------------------------------------------------------
\subsubsection*{Variations map}

\begin{diagrambox}
\begin{tikzpicture}[
    var/.style={rectangle, rounded corners=2pt, draw=patternborder, fill=patternbg,
                font=\scriptsize, align=center, text width=3.25cm, minimum height=1.3cm, inner sep=3pt}]
    \node[var, fill=answerfill, draw=accentorange, text width=5.4cm, font=\scriptsize\bfseries] (root) at (0,0)
        {rooted tree; combine children at leave, extend the parent at enter};
    \node[var] (v1) at (-5.4,-2) {\textbf{count / sum per subtree}\\colour balance\\(White-Black, 2.2.1)};
    \node[var] (v2) at (-1.8,-2) {\textbf{pieces after deletion}\\child sizes and $n - \mathrm{sz}$\\(Highest Score, 2.2.2)};
    \node[var] (v3) at (1.8,-2)  {\textbf{add to a subtree}\\mark the top, push down\\(Ki, 2.2.3)};
    \node[var] (v4) at (5.4,-2)  {\textbf{per-label counts}\\counter before / after\\(E2.1)};
    \node[var] (v5) at (-5.4,-3.6) {\textbf{greedy cuts}\\cut when complete\\(E2.2, E2.5)};
    \node[var] (v6) at (-1.8,-3.6) {\textbf{traffic on an edge}\\a function of $\mathrm{sz}$\\(E2.3)};
    \node[var] (v7) at (1.8,-3.6)  {\textbf{best path ending here}\\push-down with reset\\(E2.4)};
    \node[var] (v8) at (5.4,-3.6)  {\textbf{small summary per subtree}\\top-3 / bottom-2\\(E2.6)};
    \node[draw=patternborder, dashed, rounded corners=4pt, fit=(v1)(v4)(v5)(v8), inner sep=5pt] (grp) {};
    \draw[flow] (root.south) -- (grp.north);
    \node[note, below=2mm of grp] {all roots: rerooting (DP book Ch.~19) $\cdot$ choices: tree DP (DP book Ch.~11) $\cdot$ mixed queries: Chapter~15};
\end{tikzpicture}
\end{diagrambox}

\textbf{Twists and how to handle them}
\begin{itemize}[leftmargin=*]
    \item \textbf{Edges given as pairs $a_i < b_i$:} that does \emph{not} make $a_i$ the parent. Always root the tree yourself by a search from the root (Problem~\ref{ssec:ch2-ki}).
    \item \textbf{Parent array with $p_i < i$:} loop the labels backwards --- no adjacency lists at all (Problem~\ref{ssec:ch2-white-black}).
    \item \textbf{Parent array without that guarantee} (LeetCode's \texttt{parents[]}): build children lists and a BFS order, or peel leaves (Problem~\ref{ssec:ch2-highest-score}).
    \item \textbf{Products and sums overflow:} three pieces of $3\cdot10^4$ vertices multiply to $2.7\cdot10^{13}$; sums of $2\cdot10^5$ values up to $10^9$ need 64 bits.
    \item \textbf{Merging big summaries is expensive:} keep only what the parent can use (5 numbers in E2.6, one counter difference in E2.1) instead of whole multisets.
    \item \textbf{Updates to subtrees, answers at the end:} mark only the top vertex and push down once (Problem~\ref{ssec:ch2-ki}); the ``difference on pre-order intervals'' form of the same idea becomes a Fenwick tree when queries are interleaved (Chapter~15).
\end{itemize}

%------------------------------------------------------------
\subsubsection*{Code templates}

\begin{lstlisting}[style=cpp]
#include <bits/stdc++.h>
using namespace std;
typedef long long ll;

// (1) RECURSIVE, tree given as undirected edges: the parent parameter replaces visited[].
//     Down (enter): depth, root-path sum.  Up (leave): size, subtree sum.  Depth of recursion = height.
struct TreeRec {
    vector<vector<int>> adj;
    vector<ll> val, sub, path;
    vector<int> sz, dep;
    void dfs(int u, int p) {
        dep[u] = (p == 0) ? 0 : dep[p] + 1;          // PRE-ORDER: p is final, push down
        path[u] = (p == 0 ? 0 : path[p]) + val[u];
        sz[u] = 1; sub[u] = val[u];
        for (int v : adj[u])
            if (v != p) {
                dfs(v, u);
                sz[u] += sz[v];                       // POST-ORDER: v is final, pull up
                sub[u] += sub[v];
            }
    }
    void run(int n, const vector<pair<int, int>>& edges, const vector<ll>& w, int root) {
        adj.assign(n + 1, vector<int>());
        for (auto& e : edges) { adj[e.first].push_back(e.second); adj[e.second].push_back(e.first); }
        val = w; sub.assign(n + 1, 0); path.assign(n + 1, 0); sz.assign(n + 1, 0); dep.assign(n + 1, 0);
        dfs(root, 0);
    }
};

// (2) ITERATIVE with an order array: one BFS gives parent[] and an order with parents first.
//     Forward pass = push-down, backward pass = aggregation. Safe for n = 10^6.
struct TreeIter {
    vector<int> parent, order, sz, dep;
    vector<ll> sub, path;
    void run(int n, const vector<pair<int, int>>& edges, const vector<ll>& val, int root) {
        vector<vector<int>> adj(n + 1);
        for (auto& e : edges) { adj[e.first].push_back(e.second); adj[e.second].push_back(e.first); }
        parent.assign(n + 1, 0); order.assign(1, root);
        for (size_t i = 0; i < order.size(); i++) {
            int u = order[i];
            for (int v : adj[u]) if (v != parent[u]) { parent[v] = u; order.push_back(v); }
        }
        dep.assign(n + 1, 0); path.assign(n + 1, 0);
        for (int u : order) {                         // parents before children: push down
            int p = parent[u];
            dep[u] = (p == 0) ? 0 : dep[p] + 1;
            path[u] = (p == 0 ? 0 : path[p]) + val[u];
        }
        sz.assign(n + 1, 1); sub.assign(val.begin(), val.end());
        for (int i = n - 1; i >= 1; i--) {            // children before parents: pull up
            int u = order[i];
            sz[parent[u]] += sz[u];
            sub[parent[u]] += sub[u];
        }
    }
};

// (3) PARENT ARRAY with p[i] < i (vertex 1 is the root): the labels are already an order.
vector<int> sizesFromParents(int n, const vector<int>& p) {
    vector<int> sz(n + 1, 1);
    for (int i = n; i >= 2; i--) sz[p[i]] += sz[i];  // i's children all have labels > i
    return sz;
}

// (4) THE PART ABOVE v: deleting v leaves its child subtrees and one piece of n - sz[v] vertices.
vector<int> largestPieceAfterDeleting(int n, const TreeIter& t) {
    vector<int> worst(n + 1, 0);
    for (int v = 1; v <= n; v++) worst[v] = n - t.sz[v];          // the piece above v (0 for the root)
    for (int v = 1; v <= n; v++)
        if (t.parent[v]) worst[t.parent[v]] = max(worst[t.parent[v]], t.sz[v]);   // a piece below
    return worst;
}

int main() {
    // the tree T of the toolbox lesson; the value of vertex v is v
    int n = 8;
    vector<pair<int, int>> edges = {{1, 2}, {1, 3}, {1, 4}, {2, 5}, {2, 6}, {4, 7}, {7, 8}};
    vector<ll> val(n + 1);
    for (int v = 0; v <= n; v++) val[v] = v;
    TreeRec r; r.run(n, edges, val, 1);
    TreeIter t; t.run(n, edges, val, 1);
    vector<int> worst = largestPieceAfterDeleting(n, t);
    vector<int> p = {0, 0, 1, 1, 1, 2, 2, 4, 7};                     // the same tree as a parent array
    vector<int> sz3 = sizesFromParents(n, p);
    cout << "v  sz sub dep path worst\n";
    for (int v = 1; v <= n; v++) {
        assert(r.sz[v] == t.sz[v] && r.sub[v] == t.sub[v] && r.dep[v] == t.dep[v] && r.path[v] == t.path[v] && sz3[v] == t.sz[v]);
        cout << v << "  " << t.sz[v] << "  " << t.sub[v] << "  " << t.dep[v] << "  " << t.path[v] << "  " << worst[v] << "\n";
    }
    cout << "BFS order:";
    for (int u : t.order) cout << " " << u;
    cout << "\n";
}
\end{lstlisting}

\begin{lstlisting}[style=trace]
output: v  sz sub dep path worst          5  1  5  2  8  7
        1  8  36  0  1  3                 6  1  6  2  9  7
        2  3  13  1  3  5                 7  2  15  2  12  6
        3  1  3  1  4  7                  8  1  8  3  20  7
        4  3  19  1  5  5                 BFS order: 1 2 3 4 5 6 7 8
\end{lstlisting}

\begin{insightbox}[Which template when?]
\renewcommand{\arraystretch}{1.25}
\begin{tabularx}{\linewidth}{@{}>{\bfseries}l X@{}}
(1) \texttt{TreeRec} & short and natural; use when the height is small ($n \le 5\cdot10^4$, or a judge with a big stack) \\
(2) \texttt{TreeIter} & the default for $n$ up to $10^6$: edges in any orientation, any root, no recursion \\
(3) \texttt{sizesFromParents} & input is a parent array with $p_i < i$: no adjacency lists, no search \\
(4) \texttt{largestPieceAfterDeleting} & anything about the pieces left by deleting a vertex: child sizes and $n - \mathrm{sz}[v]$ \\
\end{tabularx}
\end{insightbox}

%------------------------------------------------------------
\subsubsection*{Common mistakes}

\begin{warningbox}[Mistake 1 --- treating the first endpoint as the parent]
\begin{lstlisting}[style=snippet]
for (each edge a b) par[b] = a;          // WRONG: "a < b" or "listed first" says nothing about rooting
// RIGHT: root the tree by a BFS/DFS from the root, skipping the parent
\end{lstlisting}
On the edges $1{-}3$, $2{-}3$ with root 1, vertex 2 is a child of 3, not the other way round. Problem~\ref{ssec:ch2-ki} shows the wrong output this produces.
\end{warningbox}

\begin{warningbox}[Mistake 2 --- using a value before it is final]
Aggregation reads the children, so it must run \emph{after} the recursive calls (or in the \emph{backward} loop). Writing \texttt{sz[u] += sz[v]} before \texttt{dfs(v, u)} adds the child's initial 0. Push-down reads the parent, so it runs \emph{before} the calls (or in the \emph{forward} loop).
\end{warningbox}

\begin{warningbox}[Mistake 3 --- forgetting the part above]
The pieces left by deleting $v$ are its child subtrees \textbf{and} $n - \mathrm{sz}[v]$. Forgetting the last one (or multiplying by it for the root, where it is $0$) is the classic bug of Problem~\ref{ssec:ch2-highest-score}.
\end{warningbox}

\begin{warningbox}[Mistake 4 --- recursion on a path of $2\cdot10^5$]
Trees in tests are often paths or ``caterpillars''. Our recursive push-down for Problem~\ref{ssec:ch2-ki} died without printing anything on a path of $2\cdot10^5$ vertices; template (2) handles it instantly.
\end{warningbox}

\begin{warningbox}[Mistake 5 --- \texttt{int} for sizes multiplied or values summed]
Sizes fit in \texttt{int}, their products do not: the first tree whose best deletion score exceeds $2^{31} - 1$ has only $n = 3872$ vertices ($1290 \cdot 1290 \cdot 1291$).
\end{warningbox}

\begin{checklistbox}[``I'm stuck'' checklist for Pattern 2]
\begin{itemize}[leftmargin=*]
    \item[$\square$] Is the answer for vertex $v$ about its \emph{subtree} (up) or about its \emph{root path} (down)? Or both?
    \item[$\square$] What does a vertex need from its children: is one number enough, or a small record (counts per colour, top 3)?
    \item[$\square$] Do I need the part above $v$? It is $n - \mathrm{sz}[v]$, or ``total minus subtree'' for sums.
    \item[$\square$] Who is the root, and is the input really oriented? Root it myself.
    \item[$\square$] How deep can the tree be? Above $5\cdot10^4$: order array.
    \item[$\square$] Do products or sums need \texttt{long long}?
    \item[$\square$] Do I need the answer for \emph{every} root, or with interleaved queries? Then this is the wrong pattern (rerooting, Chapter~15).
\end{itemize}
\end{checklistbox}
